% Ad Atticum 10.8 (ad-atticum-10-08) --- English translation
% Translated by Claude (Anthropic) from the Latin (Perseus).
% Style guide: STYLE.md.

\ciceroLetterOpener{CICERO ATTICO salutem}

\ciceroSection{1} The matter itself was warning me, and you had pointed it out, and I could see for myself, that on those subjects the interception of which would be dangerous it was time we put an end to writing between us. But since my Tullia keeps writing to me, begging me to wait for what is being done in Spain, and always adding that you are of the same view, and since I have understood this from your own letters too, I do not think it amiss to write to you what I think about the matter.

\ciceroSection{2} That plan, in my judgement, would be a prudent one if we were going to shape our course by what happens in Spain. Of necessity one of three things must come about: either, what I should most wish, that fellow is driven out of Spain; or the war is drawn out; or he, as he seems to expect, seizes the Spains. If he is driven out, how welcome or how honourable will be my arrival then with Pompey, when I judge that Curio himself will be going over to him? If the war is drawn out, what am I to wait for, and how long? That leaves the case in which, if we are beaten in Spain, we keep quiet. The opposite, in my judgement, is the truth: I think that fellow ought rather to be deserted as a victor than as one defeated, doubtful rather than confident in his fortunes. For if he wins I see slaughter, an attack on private wealth, the recall of exiles, the cancelling of debts, honours for the most shameful men, and a kingship not merely intolerable to a Roman --- intolerable to any Persian.

\ciceroSection{3} Will our outrage be able to keep silent? Will my eyes be able to bear the sight of me speaking my vote alongside Gabinius, and to bear his being called on before me? to have your client Clodius standing by, and Plaguleius the freedman of Gaius Ateius, and the rest? But why do I gather up the enemies, when I shall not be able to look upon, without pain, in the Senate House, those of my own connections whom I have defended, or to move among them without disgrace? And what if it is not even certain that I shall be allowed to do so? For his friends write to me that I have by no means satisfied him by not coming to the Senate. Are we even so to hesitate whether to sell ourselves to him --- and at peril at that --- when not even with reward would we be joined to him?

\ciceroSection{4} Then consider this too: the whole question is not being decided in Spain --- unless perhaps you think that, with the Spains lost, Pompey would throw down his arms, when his whole plan is Themistoclean. He judges that the man who holds the sea must hold the situation. For that reason it has never been his object that the Spains should be held for their own sake; the readying of a fleet has always been his oldest concern. He will sail, then, when the time is ripe, with the largest of fleets, and will come up to Italy. In which, sitting still, what shall we be? For to remain in the middle will no longer be permitted. Shall we oppose his fleets, then? What crime could be greater than that, what could equal it? What more disgraceful? Did I alone bear that man's crime against absent men, and shall I not bear it now along with Pompey and the rest of the leading men?

\ciceroSection{5} But if, with the duty now waved aside, account is to be taken of risk: from those men there is risk if I do wrong, from this one if I do right; and no plan in these troubles can be found that is free of risk, so that there is no doubt we should be running from a shameful course at peril, when we would run from it even if it were safe. \dag{}Did we not cross the sea together with Pompey? We could not at all.\dag{} The reckoning of the days makes it plain. And yet --- to confess what is the fact --- we are not even now packing up so that we can. One thing has deceived me, which perhaps it ought not to have done, but it did: I thought there would be peace. If there had been, I did not want Caesar to be angry with me when he was also Pompey's friend; for I had perceived how much of one mind they were. From fear of this I have fallen into this slowness. But I shall make up all of it if I press on; if I delay, I shall lose all.

\ciceroSection{6} And yet, my Atticus, omens too rouse me with a certain hope which is not unsure --- not those of our College of augurs from Attus Navius, but those of Plato about tyrants. For I see no way in which that fellow can stand any longer without his falling, by his own act, even with us doing nothing; seeing that, at the very height of his fortune and newly come, he has in six or seven days drawn upon himself the bitterest hatred of the very mob, needy and lost as it is; that he has so quickly stripped away the pretence of two things, mildness in the matter of Metellus, riches in the matter of the treasury. And then, what allies or what ministers has he to use? Are those men to govern provinces, are those men to govern the commonwealth, of whom not one was able to manage his own patrimony for two months?

\ciceroSection{7} It is not for me to gather up all the considerations which you with the keenest insight see for yourself; but set them before your eyes, and you will at once understand that such a kingship cannot last so much as half a year. If I am mistaken, I shall bear it --- as many of the most distinguished men, in their excellence in public affairs, have borne it --- unless you happen to think I prefer, after the manner of Sardanapallus, to die in my own little bed rather than in a Themistoclean exile. Themistocles, who, though, as Thucydides says, he was \textit{[Greek: t\=on men paront\=on di' elachist\=es boul\=es kratistos gn\=om\=on]} --- of present matters by briefest deliberation the most certain judge --- still fell into those misfortunes which he would have escaped if nothing had deceived him. And though he was, as the same writer says, the man who \textit{[Greek: to ameinon kai to cheiron en t\={o}i aphanei eti he\={o}ra malista]} --- best of all foresaw, while it was yet hidden, the better and the worse outcome --- still he did not see how he was to escape the ill-will either of the Spartans or of his own fellow citizens, nor what to promise to Artaxerxes. That night would not have been so bitter to Africanus, wisest of men; that day of Sulla's would not have been so dreadful to Gaius Marius, shrewdest of men, if nothing had deceived either of them. As for us, we are confirming ourselves in this by that omen of which I spoke; and it does not deceive us, and it will not turn out otherwise. That fellow must fall, either by his adversaries or by himself, who is indeed of all his adversaries the fiercest. I hope this will happen while we are still alive --- though it is time for us to be thinking by now of that other, everlasting life, not this brief one. But if anything should befall us sooner, it will make no very great difference to me whether I see what is being done at the time, or whether I have long foreseen that it would be. Since this is so, it must not be admitted that I obey those whom the Senate, lest the commonwealth take any harm, armed against me.

\ciceroSection{9} Everything is in your charge, which from your love for me has no need of my charging. Nor, by Hercules, can I myself find what to write; for I sit there waiting \textit{[Greek: ploudok\=on]}. Even so, nothing was ever so worth writing as that nothing has ever happened to me, out of so many kindnesses of yours, more welcome than that you have looked after my Tullia so sweetly and so attentively. She herself was much delighted by it, and I no less. Her courage, indeed, is wonderful. How she bears the public catastrophe! how she bears her domestic troubles! And what a spirit she showed at our parting! There is affection \textit{[Greek: storg\=e]}; there is the deepest fellow-feeling \textit{[Greek: synt\=exis]}. And yet she wants us to do what is right and to be well spoken of.

\ciceroSection{10} But enough of this, lest I rouse up my own sympathy \textit{[Greek: sympatheian]} too much. You, if you have anything more certain about the Spains, and anything else, while I am still here, write to me; and I too, perhaps, on parting, will send you a line, the more so because Tullia did not think you were leaving Italy at this time. With Antony we must take the same course as with Curio: that I wish to be on Malta, that I do not wish to be a party to this civil war. I should be glad to find him as easy and as good toward me as Curio. He was reported to be coming to Misenum on the sixth day before the Nones, that is to-day. But he has sent ahead to me a tiresome letter, of which this is the copy:
